% Supplementary material: moved from Section 10.7 of the main text.

\section{Open Research Directions}
\label{supp:directions}

Beyond direct empirical verification, several promising research directions emerge from this review:

\begin{description}
    \item[Multilingual \gls{slm} evaluation] Extending the consolidated performance evidence to non-English benchmarks would reveal whether the English-centric performance hierarchy reported in the literature holds across languages, particularly for models with multilingual training data (Qwen2.5, Gemma~3, BLOOM).

    \item[On-device continual learning] Developing methods for incremental adaptation of deployed \glspl{slm} to new data without full retraining would enable personalization and domain adaptation at the edge.

    \item[SLM ensembles vs.\ single \gls{llm}] Investigating whether ensembles of specialized \glspl{slm} (one for math, one for code, one for general knowledge) with a lightweight router can match single-LLM performance at lower aggregate cost.

    \item[Multimodal \glspl{slm}] Evaluating vision-language \glspl{slm} such as Phi-3-vision, Gemma~3 with vision, and Gemma~3n \citep{google2025gemma3n} on multimodal benchmarks would extend the comparison to the rapidly growing multimodal inference domain.

    \item[Hybrid architectures] The emergence of non-transformer \gls{slm} architectures (SambaY \citep{ren2025sambay}, Gated Delta Networks \citep{yang2025gated_delta}) suggests that the transformer's dominance may not extend to the edge, where linear-time inference is particularly valuable.

    \item[Federated fine-tuning] Developing privacy-preserving fine-tuning protocols that aggregate gradient updates across edge devices would enable collaborative model improvement without centralized data collection.
\end{description}
